\scanpage{125}
\mspage{40}
% archive sheet number: 124

% L01
5. \uline{Démonstration du théorème de monodromie.}

% L02-06
Nous pouvons maintenant prouver 3.2, en utilisant la suite spectrale (2.5\inserted{bis}) \inserted{de 2.10.} (pour $\nu$ variable). Les résultats du n° \deleted{\unsure{suivant}}\inserted{précédent} prouvent que les $(\underline{\Psi}_\nu^q)\inserted{_{\mathrm{tame}}}$ sont des faisceaux constructibles, donc les $E_2^{pq}$ \deleted{sont} (pour $\nu$ fixé) sont des groupes finis.

% L07-09
Donc, passant à la limite sur $\nu$, on trouve une suite spectrale, compatible avec les opérations de $I_t$ et \unsure{ce|même} de $\unsure{\pi|\Pi}$ tout entier

% L10
\[
\deleted{H^*(\overline U,\deleted{\unsure{\Lambda}}\unsure{\mathbb{Z}|Z}_\ell)^P\Longleftarrow\inserted{\unsure{E_2^{pq}=\varprojlim}}H^p(\overline X_0,(\underline{\Psi}^q)_{\mathrm{tame}})}
\tag*{\deleted{(5.1)}}
\]

% L11
\deleted{avec}
\[
\deleted{(\underline{\Psi}^q)_{\mathrm{tame}}=\varprojlim_\nu(\unsure{\overline{\underline{\Psi}}_\nu^q|\underline{\Psi}_\nu^q})_{\mathrm{tame}}}
\tag*{\deleted{(5.2)}}
\]

% L12
\[
H^*(\overline U,\mathbb{Z}_\ell)^P\Longleftarrow E_2^{pq}=\varprojlim_\nu H^p(\overline X_0,(\underline{\Psi}_\nu^q)_{\mathrm{tame}}).
\tag*{(5.1)}
\]

% L13-15
On a vu dans 4.17 que le s-groupe $I'_t\subset I_t$ opère trivialement sur les $\underline{\Psi}_\nu^q$ donc \deleted{\unsure{les}} sur les $E_2^{pq}(\nu)=H^p(\overline X_0,\underline{\Psi}_\nu^q)$, donc sur leur limite $E_2^{pq}$.

% L16-21
Donc $I'$ opère trivialement sur tous les termes de la suite spectrale (5.1), donc sur l'\deleted{\unsure{aboutissement}} gradué associé à l'aboutissement. Comme le \deleted{\unsure{filtré}} gradué est donné, \unsure{sauf extensions} à partir des facteurs, on a gagné.
